% Adapted from the frozen Phase 2 submission.
\section{Literature Review}
\label{sec:lit}
% ===================================================================

\subsection{Review Scope and Methodology}

Candidate studies were located through Google Scholar, IEEE Xplore, ScienceDirect,
SpringerLink and PubMed Central using combinations of the terms \emph{rice leaf
disease}, \emph{plant disease classification}, \emph{multiple instance
learning}, \emph{patch-based}, \emph{foundation model}, \emph{data augmentation},
\emph{disease severity} and \emph{edge deployment}. Studies were included if they
reported a quantitative evaluation on rice or on closely related foliar disease
imagery, or if they introduced a methodological component that the proposed
pipeline depends upon. Twenty-six studies are characterised in the survey tables;
of these, twenty-one were published between 2020 and 2026, with the remaining
five retained because they are the methodological origin of a technique the
project builds on. A small number of additional works are cited in the discussion
purely as methodological anchors and are not tabulated: the multiple instance
formulation \cite{dietterich1997}, attention-based instance aggregation
\cite{ilse2018}, its weakly supervised application in computational pathology
\cite{lu2021clam}, Otsu's thresholding operator \cite{otsu1979}, the paired
image-to-image translation model used for generative augmentation in the rice
literature \cite{isola2017}, and the residual and transformer backbones that
recur as baselines \cite{he2016, dosovitskiy2021}.

Tables~\ref{tab:survey1} and~\ref{tab:survey2} summarise each study by reference,
year, method, dataset, evaluation metrics, reported performance and key
limitations. Where a study reports qualitative results or metrics that are not
comparable across the table, the performance column describes the reported
outcome rather than quoting a single figure.
\clearpage



\subsection{Critical Discussion by Methodological Group}

\subsubsection{Classical image processing and handcrafted features}

The earliest rice-specific systems operated on explicitly engineered
representations. Orillo \textit{et al.} segmented lesions in MATLAB and
classified colour and texture descriptors with a back-propagation network,
reporting perfect accuracy on 134 images \cite{orillo2014}; Pugoy and Mariano
combined histogram intersection with $k$-means clustering to isolate outlier
regions \cite{pugoy2011}; Azim \textit{et al.} used a saturation threshold to
remove background and a hue threshold to isolate lesions before applying XGBoost
\cite{azim2021}; and Sethy \textit{et al.} occupy an intermediate position,
extracting CNN features but classifying them with an SVM \cite{sethy2020}. Two
observations matter for this project. First, the reported accuracies are
uninformative as generalisation estimates: a 100\,\% score on 134 images
\cite{orillo2014} and comparable figures on the 120-image UCI subset
\cite{azim2021} reflect dataset size far more than method quality. Second, and
more usefully, this group establishes the \emph{lineage} of the thresholding
operations that persist in modern pipelines. When a recent deep pipeline applies
Otsu thresholding as a preprocessing stage \cite{ali2026}, it is inheriting a
design decision from this era, with the same underlying assumption of a clean
bimodal separation between leaf and background. That assumption is testable, and
the evidence within the literature is mixed at best: one study using the same
rice dataset describes its backgrounds as essentially uniform and therefore not
requiring background subtraction \cite{tejaswini2022}, while the same operator
applied to raw images of that dataset has been reported to leave classifier
accuracy largely unchanged \cite{ali2026}.

\subsubsection{CNN transfer learning}

This is by far the largest group and it is also the most saturated. Following
Mohanty \textit{et al.} \cite{mohanty2016}, rice-specific studies have applied
essentially every mainstream ImageNet backbone to essentially the same handful of
public datasets \cite{lu2017, chen2020, ghosal2020, tejaswini2022}. The pattern
of reported results is more revealing than any individual number. On the same
Kaggle rice dataset, published accuracies span 58.15\,\% \cite{thakur2023},
74--78\,\% \cite{tejaswini2022, ali2026} and 86.98\,\% \cite{ali2026}, a spread
of nearly thirty points across differing protocols. Partitioning, preprocessing,
augmentation and metric definitions may contribute to these differences; the
reported scores do not isolate architecture. This motivates the internally
consistent baseline comparisons in the present study. The group also shares two structural
limitations. Input resolution is fixed early and never ablated, and evaluation is
almost exclusively in-domain and accuracy-based. Ghosal and Sarkar's report of
92.46\,\% training against 81\,\% test accuracy \cite{ghosal2020} is a useful
illustration of how much information a single accuracy figure conceals.

\subsubsection{Transformers and hybrid architectures}

Recent work substitutes or supplements the convolutional backbone with
self-attention \cite{alfalluji2025, vitresnet2025, globallocal2026}. The trend is
towards two-stream designs in which a transformer branch captures global context
and a convolutional branch captures local texture, fused by a learned attention
mechanism \cite{globallocal2026}. Two of these studies provide additional evaluation detail: Ramesh \textit{et al.}
report precision, recall, F1 and AUC rather than accuracy alone
\cite{vitresnet2025}, and the two-stream study includes both cross-validation and
cross-dataset testing \cite{globallocal2026}. Nevertheless, the group is subject
to diminishing returns. Accuracies of 100\,\% and 99.72\,\% on a three-class rice
task \cite{alfalluji2025} indicate that the benchmark, not the method, is the
limiting factor; notably, the same study's accuracy falls to 90.69\,\% on a
five-class task, a drop that is reported but not diagnosed. Critically, every
model in this group still performs a single forward pass over one downsampled
view of the leaf. A hierarchical transformer attends across patches \emph{within}
a $224\times224$ input, which is a different operation from attending across
regions of the original $1881\times1881$ capture; the information discarded by
the resize is not recoverable by attention applied afterwards. This distinction
is central to the gap this project targets.

\subsubsection{Generative data augmentation}

Two sub-lines exist. Adversarial approaches dominate the rice literature: the
pipeline most directly relevant to this project pairs a Pix2Pix generator with
Otsu preprocessing and an InceptionDenseNet classifier \cite{ali2026}, and more
recent work improves generator fidelity through frequency-domain and
dual-discriminator designs \cite{fhwd2025}. Diffusion-based augmentation is the
emerging alternative and has been reported to produce more realistic and more
diverse plant disease imagery than comparable GANs \cite{muhammad2023}. Three
experimental considerations follow. First, fidelity is not the objective: FID measures how
convincing an image looks, not whether adding it to the training set improves
generalisation, and the causal link between the two is assumed far more often than
it is demonstrated. Second, the rice literature contains a clear negative result
worth taking seriously --- adversarial augmentation applied on its own has been
reported to \emph{reduce} classifier accuracy relative to the raw data, with
improvement appearing only when combined with a preprocessing stage
\cite{ali2026}. Third, paired image-to-image generation such as Pix2Pix \cite{isola2017}
requires an explicit conditioning representation and paired training data. This
adds an experimental factor alongside the classifier and preprocessing. The
present study therefore uses conventional augmentation and focuses on spatial
representation and aggregation; it does not test whether synthetic augmentation
would improve this protocol.

\subsubsection{Segmentation-based preprocessing and severity quantification}

A different response to background clutter replaces
thresholding with a learned segmentation network. RCAMNet performs multi-class
segmentation with U-Net, DeepLabv3 or Detectron2 and then applies dual-path
attention for joint classification and severity staging \cite{sudhesh2025}, while
a UAV study reaches 97.20\,\% mIoU using U-Net with a ResNet-101 encoder on
multispectral imagery scored against the IRRI Standard Evaluation System
\cite{logavitool2025}. These results confirm that lesion-level supervision, when
available, is highly effective. The binding constraint is annotation. Both studies
depend on purpose-built datasets with expert pixel-level labels, and the public
rice datasets in general use provide no masks; the recently released
\textit{RiceSeg-5932} mask collection \cite{riceseg2025} addresses this for a
different rice image set but not for the datasets used here. There is also a
subtler point: segmentation removes background but does not address resolution,
since the segmented image is still resized before classification. A method that
requires no masks and preserves resolution is therefore preferable on both counts,
and this reasoning is why segmentation is treated in this project as a comparison
point rather than as the proposed solution.

\subsubsection{Patch-level, lesion-level and weakly supervised learning}

This group contains the nearest precedents for the proposed method and is, at the
same time, the thinnest group in the plant disease literature. DeChant
\textit{et al.} tile field images of maize, classify each region with a CNN,
assemble the per-region outputs into heat maps, and classify the whole image from
those maps, reaching 96.7\,\% on held-out data \cite{dechant2017}. Barbedo shows
that classifying individually cropped lesions substantially outperforms
whole-leaf classification across fourteen species \cite{barbedo2019}. Bollis
\textit{et al.} use activation maps to propose regions and then apply multiple
instance learning over them for citrus pest classification \cite{bollis2020}.
The formulation these works approach is multiple instance learning
\cite{dietterich1997}, in which a labelled \emph{bag} of unlabelled
\emph{instances} is classified as a whole; in its modern form, a learned
attention function assigns each instance a weight and the bag representation is
their weighted mean \cite{ilse2018}, a design that underpins clinical-grade
weakly supervised classification of gigapixel pathology slides
\cite{lu2021clam}. The structural correspondence to leaf imagery is close: a
high-resolution image, a diagnostic region occupying a small fraction of it, and
a label available only at whole-image level. What is missing is the combination.
DeChant \textit{et al.} require manual lesion marking to train the region
classifier and use a fixed, non-learned aggregation \cite{dechant2017}; Barbedo
requires lesions to be cropped by hand \cite{barbedo2019}; Bollis \textit{et al.}
use a two-stage saliency-then-MIL design on a pest task \cite{bollis2020}. To the
best of our knowledge, no published work applies end-to-end attention-based
instance aggregation over high-detail tiles, trained from image-level
labels alone, to rice leaf disease recognition. This is the methodological centre
of the gap identified in Section~\ref{sec:gap}.

\subsubsection{Foundation models and frozen self-supervised features}

Self-supervised vision transformers trained on web-scale unlabelled imagery
\cite{oquab2024} have recently been evaluated in agriculture with results that
are directly relevant here. Foundation features have been assessed in agriculture
under linear probing, full fine-tuning and low-rank adaptation, with the frozen
linear probe remaining competitive at a small fraction of the training cost and
foundation features outperforming other self-supervised alternatives when
distilled into a mobile backbone \cite{dinoagri2025}. More pointedly, frozen
self-supervised features have been shown to substantially reduce the
laboratory-to-field degradation that afflicts supervised backbones
\cite{nguyen2026}, and contrastive self-supervised pretraining under an explicit
leakage-controlled protocol improves robustness under domain shift
\cite{plantclr2026}. Two properties make this group attractive as a component
rather than a contribution. Because a frozen encoder is never trained on the
target data, it cannot memorise dataset-specific background cues, which is
exactly the failure mode that motivates background removal in the first place.
And because embeddings can be computed once and cached, the cost of aggregating
many regions per image becomes tractable, which is what makes a tile-based
formulation feasible under student-level compute. The gap is that all of these
studies probe features from whole downsampled images; none evaluates frozen
features as region descriptors under a weakly supervised aggregator, and none
targets rice.

\subsubsection{Datasets, evaluation protocol and generalisation}

This group supplies the evidence for the protocol gap. Xu \textit{et al.}
provide a taxonomy of plant disease datasets and argue that the apparent success
of deep learning in this field is partly illusory, identifying random splitting
as a direct cause of overestimated test performance because near-duplicate views
of one physical observation can be distributed across training and test
partitions, and observing that splitting and annotation strategy are scarcely
discussed in applied papers \cite{xu2024}. This is a field-wide reporting issue
rather than a fault of any individual study: where the ordering of augmentation
and partitioning, or the partition on which a confusion matrix was computed, is
not stated, the reader simply cannot determine what a reported figure measures.
The consequences are visible empirically. A model exceeding 90\,\% on four
benchmarks reaches 58.15\,\% on an unseen rice dataset \cite{thakur2023}, and
the laboratory-to-field drop documented for frozen-feature studies
\cite{nguyen2026, plantclr2026} is larger than the margins by which competing
architectures are separated. On the resource side, the group also supplies the
remedy. Paddy Doctor offers 16{,}225 expert-annotated field images across 13
classes with published transfer-learning baselines and an F1-based reporting
convention \cite{petchiammal2023}, and PlantDoc demonstrates the value of
deliberately in-the-wild data for breaking laboratory shortcuts
\cite{singh2020}. The availability of a large, independently collected field
dataset whose class vocabulary overlaps the primary dataset is what makes the
zero-shot cross-dataset evaluation in objective~6 practical rather than
aspirational.

\subsubsection{Lightweight models and edge deployment}

The final group addresses whether a model can run where it is needed.
MobileNet-V4 has been benchmarked against earlier mobile backbones across GPU,
CPU and Raspberry Pi 5, reporting 97.84\,\% accuracy at 88.91\,ms per inference
on the Pi \cite{edge2026}, and structured pruning combined with channel-wise
feature distillation compresses a YOLOv11 variant to 1.9\,MB while retaining
97.4\,\% mAP@50 \cite{liu2025}. These studies are notable for measuring latency
on real hardware rather than asserting deployability, a standard that much of the
rice literature does not meet --- claims that a 98\,MB classifier is suitable for
IoT or edge use are common but are rarely accompanied by an on-device
measurement. The limitation is that this group optimises for in-domain accuracy
under compression and does not test whether the compressed model retains
generalisation. For this project the group functions as a cost benchmark: any
region-based pipeline must be assessed for computational overhead against these
figures, which is the purpose of objective~7.

\subsection{Synthesis: Cross-Cutting Limitations}

Reading across the nine groups, four limitations recur with enough consistency to
be treated as properties of the field rather than of individual papers. Input
resolution is treated as a fixed preprocessing constant and is never ablated,
despite the two-order-of-magnitude information loss it entails on
high-resolution captures. Background suppression is addressed at opposite
extremes of the supervision spectrum --- hand-tuned global operators on one side
\cite{ali2026, azim2021} and densely annotated segmentation networks on the other
\cite{sudhesh2025, logavitool2025} --- with no annotation-free learned option
evaluated in between. Overall accuracy on a randomly partitioned, strongly
imbalanced dataset remains the dominant metric, so a model that never predicts a
rare but agronomically important class pays almost no measured penalty, and
repeated runs with reported variance are uncommon. Finally, cross-dataset
evaluation is rare enough that the field-readiness claims made in the motivation
sections of these papers are, in most cases, untested by their own experiments.

It is worth being explicit that the four are connected rather than independent.
Aggressive downsampling removes lesion-scale evidence, which leaves the
classifier more reliant on whatever global cues are available, including
background and acquisition artefacts; that reliance is invisible under in-domain
accuracy on a randomly split single dataset, and surfaces only under
cross-dataset evaluation, where it appears as the large degradations reported in
\cite{thakur2023, nguyen2026}. A method that preserves resolution and learns
region relevance therefore addresses the measurement problem and the modelling
problem through the same mechanism, provided the evaluation is designed to detect
the difference.

\subsection{Research Gap}
\label{sec:gap}

Three gaps follow directly from the synthesis above.

\begin{itemize}
\item \textbf{Resolution gap.} Rice disease recognition pipelines almost
universally downsample high-resolution captures to between 128 and 256 pixels
before classification, discarding the fine textural evidence that distinguishes
early-stage and visually similar lesions. No study in the reviewed literature
treats input resolution as an experimental variable, and none evaluates
resolution-preserving region aggregation as an alternative, although the
technique is standard in gigapixel medical imaging \cite{ilse2018, lu2021clam}.

\item \textbf{Supervision gap.} Suppression of diagnostically irrelevant image
content is handled either by hand-engineered global operators that assume a
bimodal leaf--background separation \cite{ali2026, azim2021, otsu1979}, or by
segmentation networks that require pixel-level masks unavailable for the public
rice datasets \cite{sudhesh2025, logavitool2025}. An annotation-free mechanism
that \emph{learns} which regions are informative from image-level labels alone
has been demonstrated in maize and citrus in partial forms
\cite{dechant2017, barbedo2019, bollis2020} but has not been evaluated for rice
in an end-to-end attention-based formulation.

\item \textbf{Protocol and generalisation gap.} Reported performance in this
field is difficult to interpret and difficult to compare, because partitioning
strategy and the ordering of augmentation relative to partitioning are often
unstated \cite{xu2024}, class-balanced metrics and run-to-run variance are often
unreported, and zero-shot evaluation on an independently collected dataset is
rare. Consequently the generalisation claims that motivate the work are largely
unsupported by the experiments presented, as the observed collapse from above
90\,\% to 58.15\,\% on an unseen rice dataset illustrates \cite{thakur2023}.
\end{itemize}

\subsection{Problem Statement}

Given high-resolution rice leaf images annotated only with an image-level
disease label, the problem is to classify the disease while (i) preserving
lesion-scale spatial detail rather than discarding it through whole-image
downsampling, (ii) suppressing diagnostically irrelevant regions through a
learned mechanism that requires no pixel-level supervision, and (iii) reporting
the result under an evaluation protocol that fixes partitioning before any
augmentation, reports class-balanced metrics with variance across repeated runs,
and, in the broader planned study, measures zero-shot transfer to an
independently collected field dataset with overlapping disease classes. The
present submission evaluates the in-domain component; external generalisation
remains untested. The proposed approach treats each leaf image as a
collection of high-detail regions at a common leaf scale and learns, from image-level labels
alone, both a class prediction and an attention distribution over regions
showing the model's learned weighting of the regions. These visual weights
are descriptive; causal faithfulness requires the separately planned tests.

\medskip

